\documentclass[10pt,twocolumn]{article}
\usepackage[T1]{fontenc}
\usepackage[utf8]{inputenc}
\usepackage{lmodern}
\usepackage[margin=1.8cm,top=2.4cm,bottom=2cm,columnsep=0.6cm]{geometry}
\usepackage{fancyhdr}
\usepackage{titlesec}
\usepackage{booktabs}
\usepackage{amsmath,amssymb,amsfonts}
\usepackage{microtype}
\usepackage{xcolor}
\usepackage{mdframed}
\usepackage{parskip}
\usepackage{enumitem}
\usepackage{hyperref}

\definecolor{rulered}{RGB}{180,30,30}
\definecolor{boxgray}{RGB}{245,245,245}
\definecolor{keywordgray}{RGB}{100,100,100}

\pagestyle{fancy}
\fancyhf{}
\fancyhead[L]{\textsc{\small Claude Mag}}
\fancyhead[C]{\small\textit{Machine Learning Research Digest}}
\fancyhead[R]{\small 2026-10-21}
\fancyfoot[C]{\small\thepage}
\renewcommand{\headrulewidth}{0.8pt}

\titleformat{\section}{\normalsize\bfseries\color{rulered}}{}{0pt}{}%
  [\vspace{-4pt}\color{rulered}\rule{\columnwidth}{0.4pt}]
\titlespacing{\section}{0pt}{8pt}{4pt}

\hypersetup{colorlinks=true,urlcolor=rulered,linkcolor=black}

\begin{document}
\setlength{\parindent}{0pt}
\setlength{\parskip}{4pt}

{\small\color{keywordgray}\textbf{Keywords:}
  robot manipulation \textbullet{}
  skill consolidation \textbullet{}
  hypernetwork \textbullet{}
  VLA models}\par\vspace{4pt}

{\LARGE\bfseries\raggedright
  NestedVLA: Learning to Consolidate and Generate Skills for
  Vision-Language-Action Models}\par\vspace{4pt}

{\large\itshape\raggedright
  A Nested Hypernetwork Generates Task-Adaptive Parameters Post Hoc,
  Outperforming Model Merging by 41 Percentage Points}\par\vspace{6pt}

{\small\textbf{By} Guodong Du \& Wanyu Lin (WanyuGroup)}\par\vspace{2pt}

{\small\color{keywordgray}arXiv:2610.23271 \textbullet{}
  NeurIPS 2026}
\par\vspace{2pt}\noindent{\color{rulered}\rule{\columnwidth}{0.6pt}}\vspace{6pt}

Teaching a robot to consolidate skills learned across heterogeneous tasks
has long been limited by cross-platform data splits, gradient interference
in joint training, and parameter conflicts in model merging. NestedVLA sidesteps
all three by learning skills independently and consolidating post hoc through
\textit{generation} rather than aggregation: a nested hypernetwork, conditioned
on task context, synthesises task-adaptive action-expert parameter deltas
on demand and composes them with a frozen base VLA. On MetaWorld, ManiSkill,
and CALVIN combined, NestedVLA outperforms the best model-merging baseline
by an average of 41 percentage points. Accepted at NeurIPS 2026.

\begin{mdframed}[backgroundcolor=boxgray,linecolor=rulered,linewidth=1pt,
    innerleftmargin=6pt,innerrightmargin=6pt,
    innertopmargin=6pt,innerbottommargin=6pt]
\textbf{\small AT A GLANCE}\par\vspace{4pt}
\begin{description}[leftmargin=0pt,labelsep=4pt,itemsep=2pt,parsep=0pt,
    font=\small\bfseries]
  \item[Problem:] Consolidating VLA skills across heterogeneous tasks
    causes gradient interference in joint training and parameter conflicts
    in model merging.
  \item[Why it matters:] A general-purpose robot must master dozens of
    skills without one degrading another; current merging methods degrade
    all.
  \item[Method:] Learn skills independently; train a nested hypernetwork
    to generate task-specific parameter deltas on demand; compose with
    frozen base VLA.
  \item[Result:] $+41$ pp average success rate over strongest model-merging
    baseline across MetaWorld, ManiSkill, and CALVIN.
  \item[Evidence:] Ablations confirm generation over aggregation is the
    key factor; adding new skills post hoc requires no global retraining.
\end{description}
\end{mdframed}

\section{The Big Picture}

A generalist robot manipulation agent must be a jack of all trades:
grasping fragile objects, pushing heavy boxes, drawing precise trajectories,
stacking irregular items. Each skill is typically learned from a dataset
collected on specific hardware, with a specific camera setup, and sometimes
on a completely different robot platform.

The dream is a single model that holds all these skills simultaneously. The
nightmare is that every approach to combining them introduces a new failure
mode: joint training causes gradient interference; model merging causes
parameter conflicts; mixture-of-experts routing requires end-to-end
co-training of all experts from the start.

NestedVLA's key insight is to separate skill \textit{learning} from skill
\textit{consolidation}, and to implement consolidation as parameter
generation rather than parameter aggregation.

\section{Why Existing Approaches Fall Short}

\textbf{Joint multi-task training:} Gradient interference arises when
different skills require conflicting parameter updates; catastrophic
forgetting prevents stable incorporation of new skills after initial
training.

\textbf{Model merging (TIES, DARE, task arithmetic):} These methods
produce a single merged checkpoint by combining parameter deltas; when
deltas conflict in sign and magnitude, the merged parameters satisfy no
task well.

\textbf{Mixture of Experts with routing:} All experts must be trained
end-to-end simultaneously; post-hoc addition of a new skill requires
full retraining to update the router.

\textbf{Adapter methods:} Per-task adapters do not share structure;
the number of parameters grows linearly with the number of tasks.

\section{Core Mechanism}

\textbf{Stage 1 --- Independent skill learning:} Each skill $k$ is trained
independently on its own data $\mathcal{D}_k$, producing a parameter delta
$\Delta\theta_k^*$ relative to a shared base VLA $\theta_0$.

\textbf{Stage 2 --- Hypernetwork-based consolidation:} A nested hypernetwork
$\mathcal{H}_\phi$ is trained to reconstruct any $\Delta\theta_k^*$ from
task context $\tau_k$ and a compressed skill token $e_k = \text{Enc}(\Delta\theta_k^*)$.
The skill tokens are stored in a compact library. At inference, the
hypernetwork generates the appropriate delta for the queried task and
composes it additively with the frozen base:
\[
\theta(\tau) = \theta_0 + \mathcal{H}_\phi(\tau,\, \{e_k\})
\]

\textbf{VLM residuals:} Task-specific shifts in text-conditioned attention
are stored separately as compressed vectors alongside the skill tokens,
capturing linguistic context the action expert alone cannot reconstruct.

\section{Technical Formulation}

For skill $k$, the effective policy parameters are
$\theta_k = \theta_0 + \Delta\theta_k^*$, and the compressed skill token
is $e_k = \text{Enc}(\Delta\theta_k^*)$.

The hypernetwork $\mathcal{H}_\phi$ is trained with a combined objective:
\[
\mathcal{L}(\phi) = \mathbb{E}_{k,(s,a)\sim\mathcal{D}_k}\!\bigl[
  \|\hat{\Delta\theta}(\tau_k) - \Delta\theta_k^*\|_F^2
  + \lambda_{\text{BC}}\,\mathcal{L}_{\text{BC}}(\theta_0
    + \hat{\Delta\theta}(\tau_k);\,s,a)\bigr]
\]
where $\hat{\Delta\theta}(\tau_k) = \mathcal{H}_\phi(\tau_k, \{e_k\})$,
$\|\cdot\|_F$ is the Frobenius norm of the delta, and $\mathcal{L}_{\text{BC}}$
is the behavioural cloning loss on skill-$k$ demonstrations.

The delta compression uses a patch tokenizer:
$\Delta\theta_k \in \mathbb{R}^{P \times Q} \mapsto e_k \in \mathbb{R}^d$,
where $d \ll PQ$. Reconstruction quality vs.\ storage is governed by
the compression ratio $r = d/(PQ)$.

\section{Learning or Inference Procedure}

\begin{enumerate}[itemsep=2pt,parsep=0pt]
  \item Train each skill independently on $\mathcal{D}_k$ to obtain
    $\Delta\theta_k^*$.
  \item Compress each delta: $e_k = \text{Enc}(\Delta\theta_k^*)$.
  \item Train $\mathcal{H}_\phi$ on all $(e_k, \tau_k, \Delta\theta_k^*)$
    triples.
  \item At inference: encode task context $\tau$; retrieve relevant $e_k$;
    run $\mathcal{H}_\phi$ to obtain $\hat{\Delta\theta}$; execute with
    $\theta_0 + \hat{\Delta\theta}$.
  \item To add a new skill: train it independently, encode its delta,
    add token to the library. No global retraining required.
\end{enumerate}

\section{What the Guarantee Says}

The Frobenius norm objective provides a formal recovery guarantee: if
$\mathcal{H}_\phi$ achieves zero training loss, then at test time it recovers
$\Delta\theta_k^*$ exactly for any task whose context $\tau_k$ was seen
during training. In practice, the reconstruction error is a measure of how
well the hypernetwork generalises to novel task phrasings; the paper shows
generalisation holds empirically to held-out phrasings not seen in training.

\section{Experimental Findings}

\begin{itemize}[itemsep=2pt,parsep=0pt]
  \item \textbf{MetaWorld (50 tasks):} NestedVLA outperforms task arithmetic,
    TIES, and DARE; gains are largest on dexterous tasks where merging
    creates severe delta conflicts.
  \item \textbf{ManiSkill:} Consistent improvements across diverse
    manipulation categories; the hypernetwork generalises to task phrasings
    not seen during consolidation training.
  \item \textbf{CALVIN:} Language-conditioned manipulation; VLM residuals
    contribute strongly here, confirming that the linguistic context they
    carry is not recoverable from the action-expert delta alone.
  \item \textbf{Average gain:} $+41$ pp over the strongest merging baseline
    across all three benchmarks combined.
\end{itemize}

\section{Ablations and Interpretation}

\textbf{Without hypernetwork (direct merge):} Returns to merging baseline
performance; the generation step is the critical ingredient.

\textbf{Without independent pretraining (jointly trained):} Hypernetwork
training becomes unstable; the fixed deltas serve as stable anchors for
the generation objective.

\textbf{Without VLM residuals:} Performance drops on CALVIN; the action
expert alone cannot reconstruct task-specific language-conditioned attention.

\textbf{Compression ratio sweep:} Aggressive compression ($r < 0.05$)
degrades fine-grained dexterous tasks while performing well on coarse
manipulation; a ratio of $r \approx 0.1$ is the practical sweet spot.

\textbf{Why generation beats aggregation:} Parameter averaging smooths
task-specific directions in parameter space; a hypernetwork preserves
them by mapping context to the appropriate structured delta directly.

\section*{Reference}

Guodong Du, Wanyu Lin.
\textbf{NestedVLA: Learning to Consolidate and Generate Skills for
Vision-Language-Action Models}.
arXiv:2610.23271, NeurIPS 2026.
\url{https://arxiv.org/abs/2610.23271}

\end{document}
